\chapter{Arquitectura de servicios}
\label{cap:AnexoB}

Este anexo describe la topología de despliegue del sistema, los servicios que lo componen, sus dependencias y las variantes de configuración disponibles. El sistema se orquesta mediante Docker Compose.

\section{Diagrama de contenedores}

La Figura~\ref{fig:c4-contenedores} muestra la vista de contenedores C4 del sistema: los cuatro componentes desplegables, sus tecnologías y los canales de comunicación entre ellos.

\begin{figure}[ht]
\centering
\resizebox{\textwidth}{!}{% C4 Level 2 — Diagrama de contenedores
% Uso: % C4 Level 2 — Diagrama de contenedores
% Uso: % C4 Level 2 — Diagrama de contenedores
% Uso: \input{figs/c4-contenedores} dentro de figure+\centering
\begin{tikzpicture}[
  every node/.style={font=\small\sffamily},
  actor/.style={
    draw=gray!55, fill=gray!10, thick,
    rectangle, rounded corners=5pt,
    minimum width=2.4cm, minimum height=2.2cm,
    align=center, text width=2.2cm
  },
  cont/.style={
    draw=aquaESI!80, fill=aquaESI!12, thick,
    rectangle, rounded corners=4pt,
    minimum width=3.2cm, minimum height=2.6cm,
    align=center, text width=3.0cm
  },
  db/.style={
    draw=turquesa!80, fill=turquesa!12, thick,
    cylinder, shape border rotate=90,
    minimum height=2.6cm, minimum width=2.8cm,
    align=center, text width=2.4cm
  },
  arr/.style={->, >=Stealth, thick, gray!60},
  biarr/.style={<->, >=Stealth, thick, gray!60},
  lbl/.style={font=\scriptsize\sffamily, fill=white,
              inner sep=1.5pt, draw=gray!30, rounded corners=2pt},
  node distance=1.0cm and 2.2cm
]

\node[actor] (usr) {
  \textbf{Codificador}\\[4pt]
  \small clínico\\[2pt]
  \scriptsize[Actor externo]
};

\node[cont, right=2.0cm of usr] (fe) {
  \textbf{Frontend}\\[6pt]
  \scriptsize Next.js 16, React 19\\TypeScript
};

\node[cont, right=2.2cm of fe] (be) {
  \textbf{Backend}\\[6pt]
  \scriptsize Elixir 1.15, Phoenix 1.8\\Commanded (CQRS/ES)
};

\node[cont, below=1.8cm of be] (ai) {
  \textbf{AI Engine}\\[6pt]
  \scriptsize Python 3.11, FastAPI\\IIC/RigoBERTa-Clinical
};

\node[db, right=2.2cm of be] (pg) {
  \textbf{PostgreSQL 15}\\[6pt]
  \scriptsize Event Store\\+ Proyecciones
};

% Relaciones
\draw[arr] (usr) -- node[lbl, above]{HTTPS} (fe);

\draw[biarr, bend left=20] (fe) to
  node[lbl, above]{\footnotesize REST API} (be);
\draw[biarr, bend right=20] (fe) to
  node[lbl, below]{\footnotesize WebSocket} (be);

\draw[biarr] (be) -- node[lbl, right]{\footnotesize HTTP POST /predict} (ai);

\draw[biarr] (be) -- node[lbl, above]{\footnotesize SQL} (pg);

% Límite del sistema
\begin{pgfonlayer}{background}
  \node[
    draw=aquaESI!60, dashed, thick, rounded corners=10pt,
    fill=aquaESI!4, inner sep=16pt,
    fit=(fe)(be)(ai)(pg),
    label={[aquaESI!80, font=\footnotesize\sffamily\bfseries]above left:%
      Sistema de Codificación CIE-10-ES}
  ] {};
\end{pgfonlayer}

\end{tikzpicture}
 dentro de figure+\centering
\begin{tikzpicture}[
  every node/.style={font=\small\sffamily},
  actor/.style={
    draw=gray!55, fill=gray!10, thick,
    rectangle, rounded corners=5pt,
    minimum width=2.4cm, minimum height=2.2cm,
    align=center, text width=2.2cm
  },
  cont/.style={
    draw=aquaESI!80, fill=aquaESI!12, thick,
    rectangle, rounded corners=4pt,
    minimum width=3.2cm, minimum height=2.6cm,
    align=center, text width=3.0cm
  },
  db/.style={
    draw=turquesa!80, fill=turquesa!12, thick,
    cylinder, shape border rotate=90,
    minimum height=2.6cm, minimum width=2.8cm,
    align=center, text width=2.4cm
  },
  arr/.style={->, >=Stealth, thick, gray!60},
  biarr/.style={<->, >=Stealth, thick, gray!60},
  lbl/.style={font=\scriptsize\sffamily, fill=white,
              inner sep=1.5pt, draw=gray!30, rounded corners=2pt},
  node distance=1.0cm and 2.2cm
]

\node[actor] (usr) {
  \textbf{Codificador}\\[4pt]
  \small clínico\\[2pt]
  \scriptsize[Actor externo]
};

\node[cont, right=2.0cm of usr] (fe) {
  \textbf{Frontend}\\[6pt]
  \scriptsize Next.js 16, React 19\\TypeScript
};

\node[cont, right=2.2cm of fe] (be) {
  \textbf{Backend}\\[6pt]
  \scriptsize Elixir 1.15, Phoenix 1.8\\Commanded (CQRS/ES)
};

\node[cont, below=1.8cm of be] (ai) {
  \textbf{AI Engine}\\[6pt]
  \scriptsize Python 3.11, FastAPI\\IIC/RigoBERTa-Clinical
};

\node[db, right=2.2cm of be] (pg) {
  \textbf{PostgreSQL 15}\\[6pt]
  \scriptsize Event Store\\+ Proyecciones
};

% Relaciones
\draw[arr] (usr) -- node[lbl, above]{HTTPS} (fe);

\draw[biarr, bend left=20] (fe) to
  node[lbl, above]{\footnotesize REST API} (be);
\draw[biarr, bend right=20] (fe) to
  node[lbl, below]{\footnotesize WebSocket} (be);

\draw[biarr] (be) -- node[lbl, right]{\footnotesize HTTP POST /predict} (ai);

\draw[biarr] (be) -- node[lbl, above]{\footnotesize SQL} (pg);

% Límite del sistema
\begin{pgfonlayer}{background}
  \node[
    draw=aquaESI!60, dashed, thick, rounded corners=10pt,
    fill=aquaESI!4, inner sep=16pt,
    fit=(fe)(be)(ai)(pg),
    label={[aquaESI!80, font=\footnotesize\sffamily\bfseries]above left:%
      Sistema de Codificación CIE-10-ES}
  ] {};
\end{pgfonlayer}

\end{tikzpicture}
 dentro de figure+\centering
\begin{tikzpicture}[
  every node/.style={font=\small\sffamily},
  actor/.style={
    draw=gray!55, fill=gray!10, thick,
    rectangle, rounded corners=5pt,
    minimum width=2.4cm, minimum height=2.2cm,
    align=center, text width=2.2cm
  },
  cont/.style={
    draw=aquaESI!80, fill=aquaESI!12, thick,
    rectangle, rounded corners=4pt,
    minimum width=3.2cm, minimum height=2.6cm,
    align=center, text width=3.0cm
  },
  db/.style={
    draw=turquesa!80, fill=turquesa!12, thick,
    cylinder, shape border rotate=90,
    minimum height=2.6cm, minimum width=2.8cm,
    align=center, text width=2.4cm
  },
  arr/.style={->, >=Stealth, thick, gray!60},
  biarr/.style={<->, >=Stealth, thick, gray!60},
  lbl/.style={font=\scriptsize\sffamily, fill=white,
              inner sep=1.5pt, draw=gray!30, rounded corners=2pt},
  node distance=1.0cm and 2.2cm
]

\node[actor] (usr) {
  \textbf{Codificador}\\[4pt]
  \small clínico\\[2pt]
  \scriptsize[Actor externo]
};

\node[cont, right=2.0cm of usr] (fe) {
  \textbf{Frontend}\\[6pt]
  \scriptsize Next.js 16, React 19\\TypeScript
};

\node[cont, right=2.2cm of fe] (be) {
  \textbf{Backend}\\[6pt]
  \scriptsize Elixir 1.15, Phoenix 1.8\\Commanded (CQRS/ES)
};

\node[cont, below=1.8cm of be] (ai) {
  \textbf{AI Engine}\\[6pt]
  \scriptsize Python 3.11, FastAPI\\IIC/RigoBERTa-Clinical
};

\node[db, right=2.2cm of be] (pg) {
  \textbf{PostgreSQL 15}\\[6pt]
  \scriptsize Event Store\\+ Proyecciones
};

% Relaciones
\draw[arr] (usr) -- node[lbl, above]{HTTPS} (fe);

\draw[biarr, bend left=20] (fe) to
  node[lbl, above]{\footnotesize REST API} (be);
\draw[biarr, bend right=20] (fe) to
  node[lbl, below]{\footnotesize WebSocket} (be);

\draw[biarr] (be) -- node[lbl, right]{\footnotesize HTTP POST /predict} (ai);

\draw[biarr] (be) -- node[lbl, above]{\footnotesize SQL} (pg);

% Límite del sistema
\begin{pgfonlayer}{background}
  \node[
    draw=aquaESI!60, dashed, thick, rounded corners=10pt,
    fill=aquaESI!4, inner sep=16pt,
    fit=(fe)(be)(ai)(pg),
    label={[aquaESI!80, font=\footnotesize\sffamily\bfseries]above left:%
      Sistema de Codificación CIE-10-ES}
  ] {};
\end{pgfonlayer}

\end{tikzpicture}
}
\caption{Diagrama de contenedores C4 del sistema CIE-10-ES.}
\label{fig:c4-contenedores}
\end{figure}

\section{Topología de contenedores}

El sistema está formado por siete servicios Docker~\cite{merkel2014docker}. La Tabla~\ref{tab:servicios} resume su función, tecnología, puerto expuesto y dependencias.

\begin{table}[h]
\centering
\caption{Servicios del sistema y sus características de despliegue.}
\label{tab:servicios}
\begin{tabular}{llrll}
\hline
\textbf{Servicio} & \textbf{Tecnología} & \textbf{Puerto} & \textbf{Depende de} & \textbf{Rol} \\
\hline
\texttt{frontend}   & Next.js 16, React 19, TypeScript & 3000  & backend         & Interfaz web \\
\texttt{backend}    & Elixir 1.15, Phoenix 1.8         & 4000  & db, ai\_engine  & Orquestador \\
\texttt{ai\_engine} & Python 3.11, FastAPI             & 8000  & —               & Inferencia IA \\
\texttt{db}         & PostgreSQL 15                    & 5432  & —               & Persistencia \\
\texttt{mailpit}    & Mailpit                          & 8025 (UI), 1025 (SMTP) & — & Email (dev) \\
\texttt{cie10\_ml\_base} & CUDA + PyTorch + Transformers & —   & —               & Imagen base ML \\
\texttt{training}   & Python, Jupyter Lab              & 8888  & cie10\_ml\_base & Entrenamiento \\
\hline
\end{tabular}
\end{table}

\section{Flujo de arranque}

El orden de arranque impuesto por las dependencias de Docker Compose es el siguiente:

\begin{enumerate}
    \item \textbf{PostgreSQL} (\texttt{db}) arranca primero; expone el puerto 5432 internamente.
    \item \textbf{AI Engine} arranca en paralelo con la base de datos; carga el modelo en memoria durante el arranque (puede tardar entre 30 y 120 segundos según el hardware).
    \item \textbf{Backend} espera a que \texttt{db} y \texttt{ai\_engine} estén disponibles; ejecuta las migraciones Ecto al arrancar y crea las tablas del \emph{event store} si no existen.
    \item \textbf{Frontend} arranca una vez el backend está activo.
    \item \textbf{Mailpit} arranca en paralelo; solo se usa en entornos de desarrollo.
\end{enumerate}

\section{Variantes de despliegue}

El repositorio incluye tres ficheros de sobreescritura de Docker Compose que adaptan la configuración al entorno de ejecución:

\begin{description}
    \item[\texttt{docker-compose.gpu.yml}] Configuración por defecto para máquinas con GPU NVIDIA. El servicio \texttt{ai\_engine} arranca con \texttt{DEVICE=cuda} y acceso al \emph{runtime} NVIDIA. Requiere el controlador NVIDIA y el toolkit de contenedores de NVIDIA instalados en el anfitrión. Tiempo de inferencia aproximado: 11 segundos por informe en tarjetas Ampere o superior.

    \item[\texttt{docker-compose.cpu.yml}] Configuración para máquinas sin GPU. Sobreescribe \texttt{DEVICE=cpu} en \texttt{ai\_engine}. El backend de atención cambia automáticamente a SDPA (\emph{Scaled Dot-Product Attention}). Tiempo de inferencia aproximado: 30 segundos por informe.

    \item[\texttt{docker-compose.mock.yml}] Sustituye el servicio \texttt{ai\_engine} por \texttt{ai\_engine\_mock}, una implementación ligera en FastAPI que devuelve predicciones ficticias sin cargar ningún modelo. Recomendado para el desarrollo del frontend y del backend cuando no se necesita inferencia real.
\end{description}

\section{Variables de entorno}

La Tabla~\ref{tab:envvars} recoge las variables de entorno configurables del sistema, con sus valores por defecto en el entorno de desarrollo.

\begin{table}[h]
\centering
\caption{Variables de entorno del sistema.}
\label{tab:envvars}
\begin{tabular}{lll}
\hline
\textbf{Variable} & \textbf{Valor por defecto} & \textbf{Consumidor} \\
\hline
\multicolumn{3}{l}{\textit{Base de datos}} \\
\texttt{POSTGRES\_USER}      & \texttt{postgres}   & db \\
\texttt{POSTGRES\_PASSWORD}  & \texttt{changeme}   & db, backend \\
\texttt{DATABASE\_URL}       & \texttt{postgresql://postgres:changeme@db:5432/cie10\_app} & backend \\
\texttt{EVENT\_STORE\_URL}   & \texttt{postgresql://postgres:changeme@db:5432/cie10\_eventstore} & backend \\
\hline
\multicolumn{3}{l}{\textit{Backend}} \\
\texttt{SECRET\_KEY\_BASE}   & (mínimo 64 bytes en prod.) & backend \\
\texttt{MIX\_ENV}            & \texttt{dev}        & backend \\
\texttt{BACKEND\_PORT}       & \texttt{4000}       & backend \\
\hline
\multicolumn{3}{l}{\textit{Frontend}} \\
\texttt{NEXT\_PUBLIC\_API\_URL} & \texttt{http://localhost:4000/api} & frontend \\
\texttt{NEXT\_PUBLIC\_WS\_URL}  & \texttt{ws://localhost:4000/socket} & frontend \\
\texttt{NODE\_ENV}           & \texttt{development} & frontend \\
\hline
\multicolumn{3}{l}{\textit{Motor de IA}} \\
\texttt{MODEL\_NAME}         & \texttt{IIC/RigoBERTa-Clinical} & ai\_engine \\
\texttt{DEVICE}              & \texttt{cuda}       & ai\_engine \\
\texttt{HF\_TOKEN}           & (token HuggingFace) & ai\_engine, training \\
\hline
\multicolumn{3}{l}{\textit{Email}} \\
\texttt{SMTP\_HOST}          & \texttt{mailpit}    & backend \\
\texttt{SMTP\_PORT}          & \texttt{1025}       & backend \\
\texttt{APP\_HOST}           & \texttt{localhost}  & backend \\
\hline
\end{tabular}
\end{table}

\section{Estructura del repositorio}

\begin{description}
    \item[\texttt{ai\_engine/}] Servicio FastAPI de inferencia. Contiene \texttt{main.py} (punto de entrada), \texttt{classifier.py} (clase \texttt{CIE10Classifier}), \texttt{train.py} (entrenamiento), \texttt{baseline\_dict.py} (\emph{fallback} por diccionario) y \texttt{model/} (artefactos del modelo entrenado).
    \item[\texttt{ai\_engine\_mock/}] Implementación simulada del motor de IA para desarrollo.
    \item[\texttt{backend/}] Aplicación Phoenix. \texttt{lib/app/} contiene la lógica de negocio (agregados, comandos, eventos, proyectores); \texttt{lib/app\_web/} contiene la capa HTTP (controladores, canales, \emph{plugs}, router).
    \item[\texttt{frontend/}] Aplicación Next.js. \texttt{src/components/} contiene los componentes de UI; \texttt{src/contexts/} los contextos de estado; \texttt{src/lib/} la configuración y utilidades de red.
    \item[\texttt{training/}] Entorno de entrenamiento aislado. \texttt{csv\_import\_scripts/} para descarga de datos; \texttt{data\_analysis/} para análisis exploratorio; \texttt{bert-classifier/} para notebooks de entrenamiento.
    \item[\texttt{tfg/}] Contiene los archivos del trabajo de fin de grado, incluyendo capítulos, anexos y figuras.
    \item[] 
\end{description}
